% vim: spl=pt
\begin{exercício}{Relação entre espectro residual e espectro pontual}{ex13}
    Seja \(T \in \bounded(\hilbert)\). Mostre que
    \begin{enumerate}[label=(\alph*)]
      \item se \(\lambda \in \sigma_r(T),\) então \(\conj{\lambda} \in \sigma_p(T^*);\) e
      \item se \(\lambda \in \sigma_p(T),\) então \(\conj{\lambda} \in \sigma_p(T^*) \cup \sigma_r(T^*).\)
    \end{enumerate}
\end{exercício}
\begin{proof}[Resolução de (a)]

\end{proof}
\begin{proof}[Resolução de (b)]

\end{proof}
